\section{A problem the synthetic families were not standing in for}
\label{sec:portfolio}

Every experiment so far draws its Hessian as $G = BB\T + nI$ for Gaussian $B$.
Section~\ref{ssec:limits} names that as a limitation and says the bias runs against
the method the paper advocates. This section makes the claim concrete rather than
leaving it as a caveat, by solving the problem the method was built for --- the
long-only portfolio --- on two real return series.

The program is~\eqref{eq:qp2} with the budget as the single equality and
non-negativity as the bounds,
\begin{equation}
  \min_{w}\; \tfrac12 w\T \Sigma w - \rho\,\mu\T w
  \quad\text{subject to}\quad
  \mathbf{1}\T w = 1, \quad w \ge 0,
  \label{eq:portfolio}
\end{equation}
with $\Sigma$ the annualised sample covariance of daily returns and $\mu$ the
sample mean. Two universes, so that nothing below is a property of one market:
the S\&P~500 at $n = \qpPfSpN{}$ assets over $T = \qpPfSpT{}$ days, and the
FTSE~100 at $n = \qpPfFtseN{}$ over $T = \qpPfFtseT{}$.

The FTSE series carries one instrument quoted flat across the whole window. Its
sample variance is zero, so $\Sigma$ is singular --- rank $86$ of $87$ --- and the
solver refuses the problem, correctly: an asset with no variance has no place in a
variance-minimising portfolio. The degenerate column is dropped and
$n = \qpPfFtseN{}$ is the result.

\subsection{Real covariances are the hard case, and in the expected direction}

Table~\ref{tab:portfolio} reports the minimum-variance corner, $\rho = 0$.

\begin{table}[ht]
\centering
\small
\caption{Long-only minimum variance on real data: time per solve and achieved KKT
residual. Sizes, conditioning and active-set size are given per universe.}
\label{tab:portfolio}
\begin{tabular}{lrr}
\toprule
Solver & ms & residual \\
\midrule
\quadprogPortfolioRows
\bottomrule
\end{tabular}
\end{table}

Three differences from the synthetic families, all of them in the same direction.

\paragraph{The conditioning is worse by orders of magnitude.} $\kappa(\Sigma)$ is
$\qpPfSpKappa$ on the S\&P~500 against the $O(n)$ of a Gaussian $BB\T + nI$. A
sample covariance of returns is close to low rank plus diagonal, which is exactly
the structure that makes it ill-conditioned, and it is the structure real data
has.

\paragraph{The active set is nearly everything.} At the minimum-variance corner
the solver holds $\qpPfSpActive{}$ of $\qpPfSpN{} + 1$ constraints active on the
S\&P~500, and the optimal portfolio holds only $\qpPfSpNames{}$ names. That is not
a quirk: minimum variance concentrates, and long-only minimum variance concentrates
onto a vertex at which almost every bound binds. A well-conditioned Hessian spreads
the optimum instead, which is why the synthetic families produce small active sets.

\paragraph{And so the exact walk is far longer.} It takes $\qpPfSpOuter{}$ outer
iterations here against the $\qpPdasBudgetOuter{}$ that the synthetic
budget-plus-bounds family produced at $n = 100$ in Section~\ref{ssec:pdas}. The
iteration count grows with the size of the set the walk must reach, and on real
data that set is large. Every statement in Section~\ref{ssec:pdas} about the
advantage of guessing the active set was therefore measured at the wrong end of
its range: on the S\&P~500 the fast path settles in $\qpPfSpRepairs{}$ repairs
against $\qpPfSpOuter{}$ outer iterations, and is $\qpPfFastSpeed\times$ faster
than the exact walk rather than the two-to-three of the synthetic case.

The comparison against the other solvers moves the same way. On the S\&P~500 the
exact walk is faster than the reference C implementation of the same algorithm ---
this is the large-$n$ side of the crossover of Section~\ref{ssec:compare} --- while
on the FTSE~100, at $n = \qpPfFtseN{}$, the reference is faster, which is the
small-$n$ side. Two points on one curve, on real data, without either being
constructed for the purpose.

\paragraph{Accuracy survives the conditioning.} The active-set method returns
$\qpPfSpResid$ on the S\&P~500 despite $\kappa(\Sigma) = \qpPfSpKappa$, while the
interior-point solver returns a residual five orders of magnitude larger at
the same requested tolerance, and larger again on the FTSE~100. The reason is the one given in
Section~\ref{ssec:compare} and it is sharper here: the optimum is a vertex with
$\qpPfSpActive{}$ active constraints, which is the worst case for a method that
approaches the boundary asymptotically and the natural case for one that
terminates on it.

\subsection{The frontier is what warm starting is for}
\label{ssec:frontier}

A frontier sweep varies $\rho$ in~\eqref{eq:portfolio} and changes nothing else.
That is precisely the family of Section~\ref{sec:sweep}: $\Sigma$, $C$ and $b$
fixed, only the linear term moving. Figure~\ref{fig:portfolio} traces
$\qpPfPoints{}$ points across the S\&P~500 frontier.

\begin{figure}[ht]
\centering
\includegraphics[width=0.86\textwidth]{graphs/quadprog_portfolio.pdf}
\caption{Left: the long-only efficient frontier, with the individual assets it
dominates. Right: the number of names held along it, on a log scale, falling from
$\qpPfNamesHi{}$ at the minimum-variance corner to $\qpPfNamesLo{}$ at the
return-seeking end, where the budget concentrates in a single asset.}
\label{fig:portfolio}
\end{figure}

Reusing the factorisation costs $\qpPfWarm$\,ms per point against
$\qpPfCold$\,ms solving each point from scratch, a factor of $\qpPfSweepSpeed$,
and the answers are identical to the cold solves because a reused point is
returned only against the certificate of Section~\ref{ssec:sweepcost}.

The interesting part is where that factor comes from, and it is not where one
would expect. The cached active set stayed optimal on only $\qpPfHits{}$ of the
$\qpPfPoints{}$ points; the other $\qpPfMisses{}$ were misses. A sweep spanning
three decades of $\rho$ moves a long way between consecutive points --- the
right-hand panel shows the holdings falling from $\qpPfNamesHi{}$ names to
$\qpPfNamesLo{}$ across it --- and the active set moves with them. Sampling the
frontier more finely would raise the hit rate directly; nothing here is tuned to
produce one.

So the saving is not mostly the hit path. It is that a \emph{miss} is repaired
rather than restarted. Proposition~\ref{prop:repair} turns a stale set into a
dual-feasible state by dropping the constraints whose multipliers have gone
negative, and the iteration resumes from there instead of from the unconstrained
minimum, keeping $J$ and $R$ throughout. A sweep that misses two points in three
is still $\qpPfSweepSpeed\times$ faster than solving each point cold, which is a
sharper statement about the machinery than a high hit rate would have been: it
says the mechanism that matters is the one that runs when the cache is
\emph{wrong}.

An earlier version of this experiment spaced $\rho$ linearly. Almost every point
then landed on the return-seeking corner, where the solution is a single asset and
does not move, and the cache hit $48$ times out of $50$. That number was real and
any conclusion drawn from it would have been worthless, because the sweep was
solving the same problem forty-nine times rather than tracing a frontier. The
geometric spacing is what makes the left-hand panel a curve, and the honest hit
rate is the one that comes with it.

\begin{remark}[What this section does and does not claim]
It claims that the method's advantages are larger on the problems it was designed
for than on the synthetic families used to test it, and that the direction of that
gap was predicted rather than discovered after the fact. It does not claim that
this is the best way to build a portfolio. Minimum variance on a raw sample
covariance is a poor estimator, and the estimation problem is a separate
literature;
here $\Sigma$ is a source of realistically shaped matrices, and nothing in this
paper depends on how it was estimated.
\end{remark}
